% © 2026 Ing. David Jaroš — CC BY-NC-ND 4.0
%
% This work is licensed under a Creative Commons Attribution-NonCommercial-NoDerivatives
% 4.0 International License (CC BY-NC-ND 4.0).
%
% License History: Earlier drafts (up to v0.3) were released under CC BY 4.0.
% From v0.4 onward, all material is released under CC BY-NC-ND 4.0 to protect
% the integrity of the theoretical work during ongoing academic development.
%
% See LICENSE.md for full license text.

% chronofactor_operator.tex
% Chronofactor as an Explicit Hilbert-Space Operator
%
% TARGET 2 of the 2026-04-30 decisive push.
% This document replaces the verbal definition in chronofactor_projection.md
% with an explicit operator acting on the winding Hilbert space.
%
% Status: DEFINED [L0/L1]
% Companion: canonical/alpha/chronofactor_projection.md (predecessor)
%            canonical/fields/biquaternion_time.tex (τ = t + iψ definition)

\ifdefined\INCLUDEMODE
\else
  \documentclass[12pt,a4paper]{article}
  \usepackage[a4paper, margin=2.5cm]{geometry}
  \usepackage{amsmath, amssymb, amsthm}
  \usepackage{mathtools}
  \usepackage{hyperref}
  \usepackage{xcolor}
  \usepackage{mdframed}
  \usepackage{booktabs}
  \usepackage{tcolorbox}

  \newtheorem{theorem}{Theorem}[section]
  \newtheorem{lemma}[theorem]{Lemma}
  \newtheorem{proposition}[theorem]{Proposition}
  \newtheorem{corollary}[theorem]{Corollary}
  \theoremstyle{definition}
  \newtheorem{definition}[theorem]{Definition}
  \theoremstyle{remark}
  \newtheorem{remark}[theorem]{Remark}

  \newmdenv[backgroundcolor=yellow!20, linecolor=orange, linewidth=1.5pt]{gapbox}
  \newmdenv[backgroundcolor=green!15, linecolor=green!60!black, linewidth=1.5pt]{resultbox}
  \newmdenv[backgroundcolor=red!8,   linecolor=red!60,   linewidth=1.5pt]{warnbox}
  \newmdenv[backgroundcolor=blue!8,  linecolor=blue!50,  linewidth=1.5pt]{insightbox}

  \title{\textbf{Chronofactor as a Hilbert-Space Operator}\\
         \large Explicit Definition: Winding Number, Translation, and Projection}
  \author{Ing.\ David Jaro\v{s}}
  \date{2026-04-30}

  \begin{document}
  \maketitle
\fi

%──────────────────────────────────────────────────────────────────────────────
\section{Purpose and Status}
\label{sec:chron:purpose}
%──────────────────────────────────────────────────────────────────────────────

\begin{tcolorbox}[colback=blue!4!white, colframe=blue!50!black,
                  title=Chronofactor — Operator Formulation]
\begin{tabular}{ll}
Goal    & Replace verbal definition with explicit Hilbert-space operator \\
Status  & DEFINED [L0] for winding operator and translation operator \\
        & DEFINED [L1] for projection map $\Pi: \operatorname{Im}\mathbb{H}\to\mathbb{C}_\tau$ \\
Input   & $\mathcal{B}=\mathbb{C}\otimes\mathbb{H}$, $\tau=t+i\psi$, $S^1_\psi$ compactification \\
Output  & Three explicit operators: $\hat{N}$, $\hat{T}_\psi(\delta)$, and $\Pi$ \\
\end{tabular}
\end{tcolorbox}

\begin{warnbox}
\textbf{No-fit rule}: No parameter in this document may be chosen to reproduce
$\alpha^{-1} = 137$ or $N_{\mathrm{eff}} = 12$.  All definitions follow from
$\mathcal{B} = \mathbb{C}\otimes\mathbb{H}$ axioms and the $S^1_\psi$ compactification.
\end{warnbox}

%──────────────────────────────────────────────────────────────────────────────
\section{The Winding Hilbert Space}
\label{sec:chron:hilbert}
%──────────────────────────────────────────────────────────────────────────────

\begin{definition}[Winding Hilbert space]
\label{def:chron:Hwind}
The \textbf{winding Hilbert space} of the UBT imaginary-time circle $S^1_\psi$
is the separable Hilbert space:
\begin{equation}
  \mathcal{H}_{\mathrm{wind}}
  = L^2(S^1_\psi, \mathbb{C})
  = \bigl\{ f: [0,2\pi R_\psi) \to \mathbb{C}
    \;\big|\; \|f\|^2 = \int_0^{2\pi R_\psi} |f(\psi)|^2 d\psi < \infty \bigr\},
  \label{eq:chron:Hwind_L2}
\end{equation}
with inner product $\langle f, g \rangle = \int_0^{2\pi R_\psi} \bar{f}(\psi)\,g(\psi)\,d\psi$.

Equivalently, in the winding-number basis, $\mathcal{H}_{\mathrm{wind}} = \ell^2(\mathbb{Z})$:
\begin{equation}
  \mathcal{H}_{\mathrm{wind}}
  = \overline{\operatorname{span}_\mathbb{C}}\bigl\{ |n\rangle \;\big|\; n \in \mathbb{Z} \bigr\},
  \qquad
  \langle m | n \rangle = \delta_{mn},
  \label{eq:chron:Hwind_ell2}
\end{equation}
where $|n\rangle$ denotes the $n$-th winding mode.
\end{definition}

\begin{definition}[Winding mode functions]
\label{def:chron:modes}
The \textbf{winding mode functions} (orthonormal basis of $L^2(S^1_\psi)$) are:
\begin{equation}
  \Phi_n(\psi)
  = \langle \psi | n \rangle
  = \frac{1}{\sqrt{2\pi R_\psi}}\, e^{in\psi/R_\psi},
  \qquad n \in \mathbb{Z},\quad \psi \in [0, 2\pi R_\psi).
  \label{eq:chron:winding_modes}
\end{equation}
These satisfy:
\begin{equation}
  \int_0^{2\pi R_\psi} \overline{\Phi_m(\psi)}\,\Phi_n(\psi)\,d\psi = \delta_{mn},
  \qquad
  \sum_{n \in \mathbb{Z}} \Phi_n(\psi)\,\overline{\Phi_n(\psi')} = \delta(\psi-\psi').
  \label{eq:chron:completeness}
\end{equation}
\end{definition}

\begin{remark}[The function $\Phi_n$ is the chronofactor]
\label{rem:chron:Phin_is_cf}
The function $\Phi_n(\psi) = (2\pi R_\psi)^{-1/2} e^{in\psi/R_\psi}$ is the
\textbf{chronofactor of winding number $n$}.  The name captures the idea that
the imaginary time $\psi$ acts as a phase variable: the winding number $n$
determines how many times the field phase winds around $S^1_\psi$ as
$\psi$ traverses the circle once.
\end{remark}

%──────────────────────────────────────────────────────────────────────────────
\section{The Three Chronofactor Operators}
\label{sec:chron:operators}
%──────────────────────────────────────────────────────────────────────────────

\subsection{Winding number operator $\hat{N}$}

\begin{definition}[Winding number operator]
\label{def:chron:Nhat}
The \textbf{winding number operator} $\hat{N}: \mathcal{H}_{\mathrm{wind}} \to \mathcal{H}_{\mathrm{wind}}$
is defined by:
\begin{equation}
  \hat{N}\,|n\rangle = n\,|n\rangle,
  \qquad n \in \mathbb{Z}.
  \label{eq:chron:Nhat}
\end{equation}
In the $L^2(S^1_\psi)$ representation:
\begin{equation}
  (\hat{N}\,f)(\psi)
  = \frac{R_\psi}{i}\,\frac{\partial f}{\partial \psi}(\psi)
  = -i R_\psi\,\partial_\psi f(\psi).
  \label{eq:chron:Nhat_diff}
\end{equation}
$\hat{N}$ is self-adjoint on its natural domain (absolutely continuous functions
on $S^1_\psi$ with $L^2$ derivative).
\end{definition}

\begin{resultbox}
$\hat{N}$ is the \textbf{U(1) charge generator} of the imaginary-time circle.
Eigenvalue $n$ is the winding charge (Kaluza-Klein momentum on $S^1_\psi$).
The spectrum is $\sigma(\hat{N}) = \mathbb{Z}$.
Status: [L0] — standard compact U(1) quantisation.
\end{resultbox}

\subsection{Imaginary-time translation operator $\hat{T}_\psi(\delta)$}

\begin{definition}[Imaginary-time translation]
\label{def:chron:Tpsi}
For $\delta \in \mathbb{R}$, the \textbf{imaginary-time translation operator} is:
\begin{equation}
  \hat{T}_\psi(\delta)
  = e^{i\delta\hat{N}/R_\psi}
  = \exp\!\bigl(i\delta\hat{N}/R_\psi\bigr)
  \;\in\; \mathcal{U}(\mathcal{H}_{\mathrm{wind}}),
  \label{eq:chron:Tpsi}
\end{equation}
where $\mathcal{U}(\mathcal{H}_{\mathrm{wind}})$ denotes the unitary operators on
$\mathcal{H}_{\mathrm{wind}}$.

Explicitly, in the winding basis:
\begin{equation}
  \hat{T}_\psi(\delta)\,|n\rangle
  = e^{in\delta/R_\psi}\,|n\rangle.
  \label{eq:chron:Tpsi_action}
\end{equation}
In the $L^2$ representation: $(\hat{T}_\psi(\delta)\,f)(\psi) = f(\psi + \delta)$.
\end{definition}

\begin{proposition}[Group structure]
\label{prop:chron:group}
The operators $\{\hat{T}_\psi(\delta) : \delta \in [0, 2\pi R_\psi)\}$ form a
unitary representation of $U(1) \cong S^1_\psi$:
\begin{align}
  \hat{T}_\psi(\delta_1)\,\hat{T}_\psi(\delta_2) &= \hat{T}_\psi(\delta_1 + \delta_2),
  \label{eq:chron:T_group1}\\
  \hat{T}_\psi(0) &= \mathbf{1},
  \label{eq:chron:T_group2}\\
  \hat{T}_\psi(\delta)^\dagger &= \hat{T}_\psi(-\delta) = \hat{T}_\psi(\delta)^{-1}.
  \label{eq:chron:T_group3}
\end{align}
\end{proposition}

\begin{proof}
Follows directly from $e^{i(\delta_1+\delta_2)\hat{N}/R_\psi} = e^{i\delta_1\hat{N}/R_\psi}e^{i\delta_2\hat{N}/R_\psi}$
(valid since $\hat{N}$ commutes with itself).
\end{proof}

\subsection{Chronofactor projection map $\Pi$}

\begin{definition}[Chronofactor projection]
\label{def:chron:Pi}
Let $\operatorname{Im}\mathbb{H} = \operatorname{span}_\mathbb{R}\{\mathbf{I}, \mathbf{J}, \mathbf{K}\}$
be the imaginary quaternion subspace ($\dim_\mathbb{R} = 3$) and
$\mathbb{C}_\tau = \operatorname{span}_\mathbb{R}\{1, i\}$ be the complex time plane
($\dim_\mathbb{R} = 2$).

The \textbf{chronofactor projection} is the $\mathbb{R}$-linear map:
\begin{equation}
  \Pi: \operatorname{Im}\mathbb{H} \to \mathbb{C}_\tau,
  \qquad
  \Pi(a\mathbf{I} + b\mathbf{J} + c\mathbf{K}) = (a + b + c)\,i,
  \label{eq:chron:Pi}
\end{equation}
which sends each Im$\mathbb{H}$ phase direction to the imaginary part of complex time.
Explicitly:
\begin{equation}
  \Pi(\mathbf{I}) = i, \quad
  \Pi(\mathbf{J}) = i, \quad
  \Pi(\mathbf{K}) = i.
  \label{eq:chron:Pi_basis}
\end{equation}
\end{definition}

\begin{proposition}[Properties of $\Pi$]
\label{prop:chron:Pi_props}
The map $\Pi$ in Definition~\ref{def:chron:Pi} satisfies:
\begin{enumerate}
  \item \textbf{Surjectivity}: $\Pi$ maps onto $i\mathbb{R} \subset \mathbb{C}_\tau$
        (the imaginary axis), covering all imaginary time directions. $\operatorname{Im}(\Pi) = i\mathbb{R}$.
  \item \textbf{Isotropic}: $\Pi$ is symmetric under permutations of $\mathbf{I},\mathbf{J},\mathbf{K}$
        — all three Im$\mathbb{H}$ directions are equivalent (consistent with the
        isotropic UBT vacuum).
  \item \textbf{Kernel}: $\ker(\Pi) = \{a\mathbf{I}+b\mathbf{J}+c\mathbf{K} : a+b+c=0\}$,
        a 2-dimensional subspace of Im$\mathbb{H}$ (the plane perpendicular to
        $\mathbf{I}+\mathbf{J}+\mathbf{K}$).  These are the ``neutral'' phase directions
        that do not couple to the imaginary time.
  \item \textbf{Dimension ratio}: The effective projection from $d=3$ to $d=2$ dimensions
        gives the ratio $\lambda_{\mathrm{proj}} = \dim_\mathbb{R}(\operatorname{Im}\mathbb{H})\big/\dim_\mathbb{R}(\mathbb{C}_\tau) = 3/2$.
\end{enumerate}
\end{proposition}

\begin{proof}
All claims follow directly from the definition~\eqref{eq:chron:Pi} by linear algebra.
For property~3: $a\mathbf{I}+b\mathbf{J}+c\mathbf{K} \in \ker(\Pi) \Leftrightarrow a+b+c=0$,
which is a codimension-1 constraint on $\mathbb{R}^3$, giving $\dim\ker(\Pi) = 2$.
\end{proof}

%──────────────────────────────────────────────────────────────────────────────
\section{The Chronofactor in the Full Biquaternion Field}
\label{sec:chron:field}
%──────────────────────────────────────────────────────────────────────────────

\subsection{Action on the UBT field $\Theta(q,\tau)$}

The UBT fundamental field $\Theta(q,\tau)$ is a biquaternion-valued function.
Under imaginary time translation $\psi \to \psi + \delta$:
\begin{equation}
  \Theta(q,\tau) \xrightarrow{\hat{T}_\psi(\delta)}
  \Theta(q, t + i(\psi+\delta))
  = \sum_{n \in \mathbb{Z}} \Theta_n(q)\,e^{in(\psi+\delta)/R_\psi}
  = \hat{T}_\psi(\delta)\,\Theta(q,\tau),
  \label{eq:chron:field_action}
\end{equation}
where $\Theta_n(q) = \frac{1}{2\pi R_\psi}\int_0^{2\pi R_\psi}\Theta(q,t+i\psi)\,e^{-in\psi/R_\psi}d\psi$
are the winding mode components.

\subsection{Winding mode decomposition}

\begin{proposition}[Winding mode spectrum]
\label{prop:chron:winding_spectrum}
The mode decomposition of the biquaternion field on $S^1_\psi$ is:
\begin{equation}
  \Theta(q,\tau)
  = \sum_{n \in \mathbb{Z}} \Theta_n(q)\,\Phi_n(\psi),
  \qquad
  \hat{N}\,\Theta = \sum_{n \in \mathbb{Z}} n\,\Theta_n(q)\,\Phi_n(\psi).
  \label{eq:chron:mode_decomp}
\end{equation}
The winding number $n$ determines the charge of the mode under $U(1)_\psi$.
Stable vacuum modes satisfy $|n|$ prime (prime-winding stability,
\texttt{canonical/appendices/appendix\_alpha\_geometry.tex \S4}).
\end{proposition}

\subsection{Effective degree-of-freedom count from the chronofactor}

\begin{theorem}[Chronofactor contribution to $N_{\mathrm{eff}}$]
\label{thm:chron:Neff}
For the biquaternion field $\Theta \in \mathcal{B} = \mathbb{C}\otimes\mathbb{H}$
on the imaginary-time circle $S^1_\psi$, the chronofactor operators $\hat{N}$ and
$\hat{T}_\psi$ generate a factor of $N_{\mathrm{chron}} = 4$ in the effective
mode count:
\begin{equation}
  N_{\mathrm{chron}}
  = N_{\mathrm{helicity}} \times N_{\mathrm{charge}}
  = 2 \times 2
  = 4,
  \label{eq:chron:Nchron}
\end{equation}
where:
\begin{itemize}
  \item $N_{\mathrm{helicity}} = 2$: left/right helicity from $\hat{T}_\psi$ acting on
        the holomorphic and antiholomorphic sectors of the field (equivalently:
        clockwise/counterclockwise winding modes $\Phi_n$ and $\Phi_{-n}$);
  \item $N_{\mathrm{charge}} = 2$: particle/antiparticle from the charge-conjugation
        automorphism $z \to \bar{z}$ of $\mathbb{C}$ in $\mathcal{B}=\mathbb{C}\otimes\mathbb{H}$.
\end{itemize}
Together with $N_{\mathrm{phases}} = \dim_\mathbb{R}(\operatorname{Im}\mathbb{H}) = 3$:
\begin{equation}
  N_{\mathrm{eff}}
  = N_{\mathrm{phases}} \times N_{\mathrm{chron}}
  = 3 \times 4
  = 12.
  \label{eq:chron:Neff12}
\end{equation}
\end{theorem}

\begin{proof}
$N_{\mathrm{phases}} = 3$ is the real dimension of Im$\mathbb{H}$, an
algebraic fact [L0].

$N_{\mathrm{helicity}} = 2$: the operator $\hat{T}_\psi(\delta)$ acts on
$\mathcal{H}_{\mathrm{wind}} = L^2(S^1_\psi)$.  The eigenspace decomposition of
$\hat{N}$ separates positive-$n$ (clockwise, holomorphic) and negative-$n$
(counterclockwise, antiholomorphic) sectors, each closed under the dynamics.
This gives two independent helicity sectors.

$N_{\mathrm{charge}} = 2$: the complex conjugation $\sigma: z\mapsto\bar{z}$ on
$\mathbb{C}$ in $\mathcal{B}=\mathbb{C}\otimes\mathbb{H}$ commutes with the dynamics
and exchanges $\Theta_n \leftrightarrow \bar\Theta_{-n}$.  This is the
antiparticle/particle symmetry.  The charge operator $\hat{Q}$ defined by
$\hat{Q}\,\Theta_n = \mathrm{sign}(n)\,\Theta_n$ has eigenvalues $\pm1$,
giving two charge sectors.

All three factors are independent and multiply: $N_{\mathrm{eff}} = 3 \times 2 \times 2 = 12$.
\end{proof}

%──────────────────────────────────────────────────────────────────────────────
\section{Operator Summary}
\label{sec:chron:summary}
%──────────────────────────────────────────────────────────────────────────────

\begin{center}
\renewcommand{\arraystretch}{1.4}
\begin{tabular}{llll}
\toprule
Operator & Definition & Action & Status \\
\midrule
$\hat{N}$ & $-iR_\psi\,\partial_\psi$ & $|n\rangle \mapsto n\,|n\rangle$ & [L0] \\
$\hat{T}_\psi(\delta)$ & $e^{i\delta\hat{N}/R_\psi}$ & $|n\rangle \mapsto e^{in\delta/R_\psi}|n\rangle$ & [L0] \\
$\Phi_n(\psi)$ & $(2\pi R_\psi)^{-1/2}e^{in\psi/R_\psi}$ & Eigenfunctions of $\hat{N}$ & [L0] \\
$\Pi: \operatorname{Im}\mathbb{H}\to\mathbb{C}_\tau$ & $a\mathbf{I}+b\mathbf{J}+c\mathbf{K}\mapsto(a+b+c)\,i$ & $\lambda_{\mathrm{proj}}=3/2$ & [L1] \\
$N_{\mathrm{chron}}$ & $N_h\times N_c = 2\times 2 = 4$ & Factor in $N_{\mathrm{eff}}$ & [L0] \\
$N_{\mathrm{eff}}$ & $N_\phi\times N_{\mathrm{chron}} = 3\times 4$ & $= 12$ charged modes & [L0] \\
\bottomrule
\end{tabular}
\end{center}

\begin{resultbox}
\textbf{Target 2 result}: The chronofactor is formalized as:
\begin{enumerate}
  \item A \emph{Hilbert space} $\mathcal{H}_{\mathrm{wind}} = \ell^2(\mathbb{Z})$.
  \item A \emph{self-adjoint operator} $\hat{N}$ (winding number / imaginary-time momentum).
  \item A \emph{unitary group} $\{\hat{T}_\psi(\delta)\}$ (imaginary-time translations, $U(1)$ representation).
  \item A \emph{linear projection} $\Pi: \operatorname{Im}\mathbb{H} \to \mathbb{C}_\tau$
        with projection ratio $\lambda = 3/2$.
  \item A \emph{mode count} $N_{\mathrm{chron}} = 4$ contributing to $N_{\mathrm{eff}} = 12$.
\end{enumerate}
All five objects are defined by explicit equations above with zero free parameters.
\end{resultbox}

\paragraph{Open item.}
The \emph{projection Jacobian} — the exact measure-change under $\Pi$, needed to
prove that the 3/2 factor enters $B_{\mathrm{base}}$ rigorously — remains
CONJECTURAL.  A rigorous derivation requires:
(i) a metric on the phase space of Im$\mathbb{H}$ modes,
(ii) a precise definition of mode density on $S^1_\psi$,
(iii) the Jacobian of $\Pi$.
These are addressed in \texttt{canonical/geometry/Rpsi\_dynamical\_fix.tex}
via the heat kernel / Poisson duality argument.

\ifdefined\INCLUDEMODE
\else
  \end{document}
\fi
